\subsection{Sprint 1}

Após a retrospectiva da Sprint 0, nossa equipe recebeu a autorização para iniciar a próxima etapa do projeto. O time foi dividido em três subgrupos: um para o front-end, outro para o back-end e o último para a blockchain. Acabei ficando no grupo do back-end, que era particularmente a área de maior interesse para mim. Nesta AGES, queria entender a fundo como a aplicação back-end funcionava e estando nesse sub-grupo seria uma grande chance de compreender este assunto de uma vez por todas.

A primeira tarefa que fiz ao começo da sprint foi a continuação da modelagem do banco de dados com o meu colega, e particularmente eu achei a tarefa um pouco complicada uma vez que as especificações das User Stories estavam muito volateis no início da sprint e criamos diversas modelagens que no final foram descartadas por conta deste problema. Basicamente a estrutura do banco de dados contou com quatro entidades principais, Empresa, Usuário, Certificado e Pemissão. Muito foi discutido em relação as entidades usuário e principalmente à permissão onde ficamos em um dilema se realmente era necessário a implementação de uma tabela para as permissões e que no final, decidimos implementar a tabela para facilitar a implementação de querys. Com a criação do modelo conceitual do banco de dados começamos a projetar o modelo lógico do mesmo mas não foi terminado nesta sprint.

Além das responsabilidades relacionadas ao banco de dados, assumi algumas tarefas envolvendo métodos HTTP, nos quais já tinha alguma familiaridade advindas do projeto anterior. Durante esta sprint, concluí duas tarefas do back-end. A primeira envolveu o método PATCH e tinha como objetivo a rejeição de uma companhia que se cadastrou no sistema. Esta tarefa foi vital para mim, pois me permitiu familiarizar-me definitivamente com o back-end da aplicação. Durante sua implementação, pude compreender detalhadamente as APIs do projeto, desde o uso do Lombok até a criação de tabelas por meio do JPA. Além disso, aprendi sobre a arquitetura do projeto e suas entidades, como DTO, controller, service, entre outros. Em relação a segunda task, usando o método POST, tive que criar, em conjunto com um colega do back-end, a criação de funcionários de uma empresa, e inclusive, foi por conta desta task que mudamos o banco de dados para ter apenas uma entidade de usuários. Esta task foi um grande desafio para mim pois ao contrário da task anterior, foi necessário criar muitas estruturas do zero e o processo de debugar foi um pouco doloroso. Apesar dos percalços, com a ajuda de ferramentas como o Postman, fui capaz de concertar os problemas e terminar a task antes do code-freezing.


A sprint chegou ao fim com nossa terceira reunião com o stakeholder, na qual apresentamos todo o trabalho realizado ao longo do período. Embora não tenhamos conseguido concluir todas as tarefas planejadas para a primeira sprint, alcançamos cerca de 4/5 delas, o que deixou nosso stakeholder satisfeito. Ele demonstrou um pouco de receio sobre o uso de uma blockchain pública, mas o time deu garantias de segurança, o que o tranquilizou e permitiu que prosseguíssemos. Além das discussões com o stakeholder, conseguimos mostrar uma demonstração do progresso do projeto, e para nossa felicidade, ele também aprovou.

Apesar das dificuldades enfrentadas anteriormente, considero que esta sprint foi a que mais aprendi e produzi em comparação com a AGES passada. No entanto, sou extremamente grato pela experiência adquirida em meu primeiro projeto, pois ele foi fundamental para a minha produtividade e aprendizado nesta sprint. Os desafios anteriores me proporcionaram uma base sólida para lidar com as demandas atuais de forma mais eficiente e eficaz, aumentando minha produtividade no projeto.